% Chapter 6: Conclusions and Discussion — Handbook 8.2.7
% ~1500 words
% Critical overview, LSEPI (mandatory), future work, reflections.

\section{Conclusions and Discussion}
\label{sec:conclusions}

\subsection{Objectives Revisited}
\label{sec:conc-objectives}

\begin{enumerate}
  \item \textbf{Listings and purchases: achieved.}  The system implemented
    the full auction lifecycle: listing creation with AES-256-GCM content
    encryption, sealed-bid submission with SHA-256 commitments, deadline-based
    expiry, and decryption key delivery to the verified winner.  This is
    demonstrated by the integration test suite (232 passing tests) and three
    end-to-end tests on the playground devnet.

  \item \textbf{Sealed-bid via MPC: achieved.}  The MASCOT and Shamir
    protocols both computed the correct winner without revealing non-winning
    bid values.  The 46-line \texttt{auction\_n.mpc} program ran correctly
    for party counts from 2 to 20 (Shamir) or 2 to 10 (MASCOT, limited by
    Veilid tunnel throughput).

  \item \textbf{MPC-gated decryption: achieved (modified design).}  The
    PDD proposed 1-out-of-$N$ oblivious transfer for decryption key delivery.
    The implemented design used a simpler commitment-and-challenge scheme:
    the winner proved knowledge of their committed bid, and the seller
    released the key upon verification.  As discussed in
    Section~\ref{sec:method-verification}, this provided equivalent privacy
    because the seller already learns the winning bid as MPC party~0.

  \item \textbf{Settlement on Monero testnet: not implemented.}  This
    objective was descoped as it was not inline with the research goals of
    the project, being more suited to a production system.  The time
    originally allocated was spent on MPC tunnel reliability and the MASCOT
    vs.\ Shamir benchmark comparison.
\end{enumerate}
9
\subsection{Discussion}
\label{sec:conc-discussion}

The central question was whether MPC can be performed over a
privacy-preserving peer-to-peer overlay at all.  The answer is yes: MASCOT
and Shamir both produced correct auction results when tunnelled through
Veilid's onion-routed private routes, for party counts from 2 to 20.  No
trusted intermediary was required, no participant's IP address was exposed,
and non-winning bid values were never revealed to anyone.

\textbf{MPC over a privacy overlay is feasible.}  The Danish sugar-beet
auction~\cite{bogetoft2009_mpc_live} demonstrated MPC for auctions, but over
direct TCP among a small set of trusted servers.  This project showed that the
same class of computation works in a fully decentralised, anonymous setting.
The MPC tunnel sustained the hundreds of megabytes of traffic that MASCOT
requires (109\,MB at $N{=}3$, 6.2\,GB at $N{=}10$) and delivered correct
results at 100\% success rate up to $N{=}12$ on the 40-node devnet.

\textbf{The cost of privacy is quantifiable.}  Veilid's onion routing added
$15{-}19\times$ overhead to MASCOT computation time and $10{-}30\times$ to
Shamir, attributable to ${\sim}$18\,ms per relay hop.  This is the
fundamental cost of routing each MPC round through multiple encrypted hops
rather than direct TCP, and it is the price paid for hiding participants'
network identities.

\textbf{Route management overhead is separable and fixable.}  Route exchange
dominated total auction time (over 70\% of wall-clock duration), but this is
an artefact of Veilid's current single-active-route constraint, which forces
routes to be created and exchanged fresh for every auction.  With persistent
routes (planned for Veilid~0.6.0~\cite{rioux2026_routing_overhaul}), this
cost would be amortised to near zero, reducing a 100\,s auction to
${\sim}$15\,s.

\textbf{The security model involves real tradeoffs.}  MASCOT provides stronger
guarantees (dishonest majority) but at a bandwidth cost that limits practical
party counts.  Shamir is dramatically cheaper but relies on an honest-majority
assumption that may not hold in an open marketplace.  For small auctions
($N{=}3$), MASCOT's stronger security is particularly valuable because a
single corrupt party under Shamir breaks all privacy.

\subsection{Legal, Social, Ethical and Professional Issues}
\label{sec:conc-lsepi}

\textbf{Dual-use risk.}  A private marketplace could in principle facilitate
illicit exchange.  This risk is mitigated by design: the system operates only
on a local devnet, no deployment scripts target the public Veilid network,
no real funds are handled, and no mainnet or testnet cryptocurrency wallets
are integrated.  The project is categorised as a Category~1 prototype not
intended for distribution (PDD Section~1.6).

\textbf{Privacy vs.\ accountability.}  Sealed bids protect bidders but
complicate dispute resolution.  In a production deployment, a reputation
system or escrow mechanism would be needed to balance privacy with
accountability.

\textbf{Licensing compliance.}  Veilid is licensed under MPL-2.0 (file-level
copyleft); all modifications to Veilid files are published under MPL.  MP-SPDZ
is BSD-3-Clause (permissive).  All other dependencies are MIT or Apache-2.0.
No licensing conflicts exist.

\textbf{Professional conduct.}  No human participants were involved.  All
testing used automated processes on local infrastructure.  The BCS Code of
Conduct was followed: the work does not endanger individuals, and the
dual-use risks are acknowledged and mitigated.

\textbf{Statement on illicit uses.}  The author explicitly condemns any
attempt to adapt this research for illicit trade, exploitation, or other
illegal activity.  The work is motivated by technical research questions:
privacy-preserving auction protocols, MPC integration, and overlay
networking,rather than any intent to facilitate wrongdoing.

\subsection{Future Work}
\label{sec:conc-future}

\begin{itemize}
  \item \textbf{Threshold decryption.}  Share the decryption key among MPC
    parties using threshold cryptography, removing the seller-as-single-point-of-trust
    assumption.  This is the highest-priority improvement.
  \item \textbf{Vickrey (second-price) auctions.}  The winner pays the
    second-highest bid rather than their own.  This would require the MPC
    program to output the second-highest value and changes to the verification
    scheme, since the winner would need to prove they bid \emph{at least} the
    second price rather than proving their exact bid.
  \item \textbf{Monero testnet settlement.}  The originally planned payment
    layer would complete the end-to-end flow from bidding to settlement.
  \item \textbf{Anti-replay nonces.}  Add explicit nonces to
    \texttt{DecryptionHashTransfer} messages to prevent replay attacks
    across auctions.
  \item \textbf{Formal verification.}  Model the protocol in a formal
    verification language to mechanically prove the privacy properties
    claimed in the security analysis.
  \item \textbf{Public network evaluation.}  Benchmark on the public Veilid
    network to understand real-world route latency and anonymity set size.
  \item \textbf{Adaptive protocol selection.}  Automatically choose between
    dishonest-majority (MASCOT) and honest-majority (Shamir) based on party
    count: at small $N$, corrupting a majority is realistic and MASCOT's
    stronger guarantees justify the cost; as $N$ grows, the difficulty of
    corruption increases and the system could safely switch to Shamir for
    better scalability.
\end{itemize}

\subsection{Reflections}
\label{sec:conc-reflections}

The project's primary contribution is a proof of feasibility: multi-party
computation can be performed over a privacy-preserving overlay network, and
the cost of doing so can be measured.  Before this work, no empirical data
existed on the overhead of tunnelling MPC through onion-routed messaging.
The benchmarks confirmed that the honest-vs-dishonest-majority choice has
first-order effects on both performance and reliability, not just theoretical
security. This is consistent with findings from Araki et
al.~\cite{araki2017_honest_majority}, who demonstrated that honest-majority
protocols scale orders of magnitude better than dishonest-majority
alternatives.  This suggests a pragmatic hybrid strategy: at small party counts, where
corrupting a majority is realistic, dishonest-majority protocols like MASCOT
are worth their cost.  As the party count grows, the difficulty of
corrupting a majority increases, and the performance advantage of
honest-majority protocols becomes compelling.  At some threshold, the
threat model can safely be relaxed in exchange for dramatically better
scalability.

The most important distinction that emerged is between \emph{fundamental}
and \emph{accidental} overhead.  The MPC tunnel's 15--19$\times$ slowdown
is fundamental: it is the cost of onion routing each round through encrypted
relay hops, and it buys real privacy.  The route exchange overhead (50--90\,s
per auction) is accidental: it results from Veilid's single-active-route
constraint forcing route teardown and recreation for every auction.  With
persistent routes, this cost disappears entirely, and the system becomes
practical for interactive use (${\sim}$15\,s per auction at $N{=}3$).

The most challenging aspect was Veilid's private route lifecycle.  Silent
route death, the single-active-route constraint, and the undocumented
self-import requirement each caused multi-day debugging sessions.  The
resulting workarounds (readiness barrier, reactive refresh, DHT-backed
route exchange) are correct but fragile.  These are framework maturity
issues rather than fundamental limitations of the approach.

If starting again, the author would begin MPC integration earlier (Milestone~1
rather than Milestone~2) and spend less time on the UI.  The headless binary
and benchmark harness proved far more valuable for demonstrating and measuring
the system than the desktop interface.  The author would also invest more time
in threat modelling and, given sufficient schedule, attempt stretch goals such
as implementing malicious nodes to generate real adversarial test data.
